\chapter{Selected hints and solutions}

\section*{Chapter 1}
For the at-least-one-ace problem, count the complement: $\binom{52}{5}-\binom{48}{5}$. For divisibility by 2, 3, or 5, inclusion--exclusion gives
\[
500+333+200-166-100-66+33=734.
\]
The four-person selection counts from nine students are
$\binom94=126$, $(9)_4=3024$, and $4\binom94=504$.
The repeated-letter count for $\texttt{LEVEL}$ is $5!/(2!2!)=30$.
The coefficient of $x^2$ in $(1+3x)^5$ is $\binom52 3^2=90$.
The length-five strings using all three symbols number
$3^5-3\cdot2^5+3=150$.
For Exercise 3, choose the three letter positions, then fill all positions:
$\binom83 26^3 10^5$.
For Exercises 13--16:
\begin{enumerate}[start=13]
 \item A specified string with three successes has probability
 $(1/3)^3(2/3)^2=4/243$. There are $\binom53=10$ such strings,
 giving total probability $40/243$.
 \item The ordered count is $3^4=81$. The unordered count is
 $\binom{4+3-1}{4}=15$. A selection $AAAA$ has one ordering,
 whereas $AAAB$ has four; the multiplicity is not constant.
 \item There are 11 letters, with four I's, four S's, two P's,
 and one M. The count is $11!/(4!4!2!)=34650$.
 \item Count by subset size to obtain $\sum_k k\binom nk$.
 Alternatively, choose the distinguished member in $n$ ways and
 any subset of the remaining $n-1$ members in $2^{n-1}$ ways.
\end{enumerate}

\section*{Chapter 2}
For the maximum of two dice to equal 4, both dice must be at most 4 but not both at most 3, giving $(4^2-3^2)/36=7/36$. For two uniform points $X,Y$, the band $|X-Y|<.1$ has area $1-2(.9)^2/2=.19$.
For the weighted four-outcome space, the probability of an even outcome
is $(2+4)/10=3/5$, not $2/4$.
On a 30-day calendar, the birthday collision probability is about
$.41356$ for six people and $.53084$ for seven; thus the first size above
$1/2$ is seven. For four people on five days, there are
$5\cdot4\cdot3\cdot2=120$ distinct-birthday lists among $5^4=625$ lists,
so the collision probability is $1-120/625=101/125$.
The quantity $\binom n2/d$ sums probabilities of overlapping events.
It is an upper bound, not necessarily a probability.
For Exercises 12--16:
\begin{enumerate}[start=12]
 \item There are $5\cdot4=20$ equally likely ordered pairs of distinct
 labeled balls. Exactly $3\cdot2=6$ have two reds, giving $3/10$.
 Unordered pairs give $\binom32/\binom52=3/10$ as well.
 \item Choose the fixed pair and derange the remaining three:
 $\binom52D_3/5!=10\cdot2/120=1/6$.
 \item $\P(A_1)=1/4$ and $\P(A_1\cap A_2)=2!/4!=1/12$.
 The product of the single probabilities is $1/16$, so the events
 are dependent.
 \item The preceding recurrence gives
 $(-1)^n=(n-1)D_{n-2}-D_{n-1}$.
 Substitute into $D_n=nD_{n-1}+(-1)^n$ to get the claimed formula.
 Then $D_7=6(265+44)=1854$.
 \item The probability is $\sum_{j=1}^{\infty}4^{-j}
 =(1/4)/(1-1/4)=1/3$.
\end{enumerate}

\section*{Chapter 3}
For the coin problem,
\[
\P(D\mid HHH)=\frac{1\cdot .2}{1\cdot .2+(1/2)^3\cdot .8}=\frac23.
\]
A pairwise-not-mutual example: flip two fair coins; let $A$ be first-heads, $B$ second-heads, and $C$ equal-results.

\section*{Chapter 4}
Exercises 1--18 use discrete variables. Exercises 19--20 belong to
the later continuous section.
\begin{enumerate}
 \item The masses at $0,1,2$ are $1/4,1/2,1/4$. The CDF is $0$
 for $x<0$, $1/4$ for $0\le x<1$, $3/4$ for $1\le x<2$, and
 $1$ for $x\ge2$.
 \item $\E[X]=(1+4+9+16)/10=3$ and
 $\E[X^2]=(1+8+27+64)/10=10$, so $\Var(X)=1$.
 Also $\P(X\ge3)=(3+4)/10=7/10$.
 The equally weighted average $5/2$ uses the wrong probabilities.
 \item Write $m=\E[X]$ and expand:
 $\E[(X-m)^2]=\E[X^2]-2m\E[X]+m^2=\E[X^2]-m^2$.
 \item Let $I_i$ indicate that $i$ is fixed and write $N=\sum I_i$.
 Since $\P(I_i=1)=1/n$, linearity gives $\E[N]=1$ for $n\ge1$.
 No independence is used.
 \item One example pays $999$ with probability $1/1000$ and $-1$
 otherwise. Its mean is zero and its variance is
 $999^2/1000+999/1000=999$. A single play loses one unit with
 probability $999/1000$; a zero mean does not make the payoff certain
 or the chance of a loss small.
 \item Chebyshev gives at most $1/9$. For a fair die the mean is
 $7/2$, the standard deviation is $\sqrt{35/12}$, and every value
 is within $5/2$ of the mean. Since $3\sqrt{35/12}>5/2$, the exact
 probability is zero. The bound is valid but loose.
 \item Equal masses at $-1,1$ have mean zero and variance one.
 Another example has masses $1/4,1/2,1/4$ at
 $-\sqrt2,0,\sqrt2$; its mean is zero and second moment is also one.
 \item The event consists of
 $(4,6),(5,5),(6,4),(5,6),(6,5),(6,6)$, so its probability is $6/36=1/6$.
 \item The masses at $0,1,2$ are $4/9,4/9,1/9$. The CDF is
 $0$ below zero, $4/9$ on $[0,1)$, $8/9$ on $[1,2)$, and $1$
 on $[2,\infty)$. The first two moments are $2/3$ and $8/9$,
 giving variance $8/9-4/9=4/9$. Also $\P(X\ge1)=5/9$.
 \item The values of $Y$ at $X=0,1,2,3$ are $1,0,1,4$.
 Its masses at $0,1,4$ are $3/8,1/2,1/8$, giving mean
 $0(3/8)+1(1/2)+4(1/8)=1$.
 \item The payoff probabilities are $1/2,1/3,1/6$.
 The mean is $-1+1/3+2/3=0$, and the second moment is
 $2+1/3+8/3=5$, so the variance is five.
 The positive-payoff probability is $1/3+1/6=1/2$.
 After a fee of one, the net values are $-3,0,3$: the mean is $-1$,
 variance remains five, and the probability of a positive net payoff
 becomes $1/6$.
 \item At $r=0,1,2,3,4$, the permutation counts are $9,8,6,0,1$,
 respectively, out of $24$. Thus $\E[N]=(8+12+4)/24=1$ and
 $\E[N^2]=(8+24+16)/24=2$, giving variance one.
 Fixing three of four objects forces the fourth to be fixed as well.
 \item There are $\binom{23}{2}=253$ pairs, each matching with
 probability $1/365$, so the mean is $253/365\approx.69315$.
 Multiple pairs can match in one outcome. Markov at threshold one
 gives $\P(C\ge1)\le253/365$, whereas the exact probability is about $.50730$.
 \item The new mean is $2(7/2)-3=4$, variance is
 $4(35/12)=35/3$, and standard deviation is $\sqrt{35/3}$.
 \item Each of $I,J=1-I$ has variance $1/4$, but $I+J=1$
 has variance zero. With an independent second indicator, the
 sum instead has variance $1/4+1/4=1/2$.
 \item The sum has mean seven and variance $35/6$.
 The difference has mean zero and variance $35/6$:
 negation changes the mean's sign but not the variance.
 The average has mean $7/2$ and variance $(35/6)/4=35/24$.
 \item For $n\ge2$, expand
 $N(N-1)=\sum_{i\ne j}I_iI_j$. Each of the $n(n-1)$ ordered
 terms has expectation $1/[n(n-1)]$, so $\E[N(N-1)]=1$.
 Adding $\E[N]=1$ gives $\E[N^2]=2$ and variance one.
 For $n=1$, $N=1$ is constant and its variance is zero.
 \item Write $X-a=(X-m)+(m-a)$, where $m=\E[X]$.
 The cross term has expectation $2(m-a)\E[X-m]=0$.
 Thus $\E[(X-a)^2]=\Var(X)+(m-a)^2$, uniquely minimized by $a=m$.
 \item Normalization gives $c\int_0^2x\,dx=2c=1$, so $c=1/2$.
 The mean is $\int_0^2x^2/2\,dx=4/3$ and the second moment is
 $\int_0^2x^3/2\,dx=2$. The variance is $2-16/9=2/9$.
 \item The density is $3x^2$ on $[0,1]$, zero elsewhere.
 The tail probability is $1-(1/2)^3=7/8$.
 The first two moments are $3/4$ and $3/5$,
 hence the variance is $3/5-9/16=3/80$.
\end{enumerate}

\section*{Chapter 5}
Exercises 1--14 concern discrete families and coupon collection; Exercises 15--17 use
continuous distributions.
\begin{enumerate}
 \item The complement gives
 $1-.7^5-5(.3)(.7)^4=.47178$.
 The mean is $5(.3)=1.5$ and the variance is $5(.3)(.7)=1.05$.
 \item For $0<p<1$ and nonnegative integers $s,t$, divide
 $(1-p)^{s+t}$ by $(1-p)^s$ to get $(1-p)^t$.
 The denominator is positive. If $p=1$, only $s=0$ is admissible;
 conditioning on an event of probability zero is not defined.
 \item Here $np=2$. The exact probability is
 $.998^{1000}\approx.135064522$, compared with
 $e^{-2}\approx.135335283$ for a Poisson variable of mean two.
 The absolute error is about $.000270761$.
 \item The mean is $9/2$, variance $(8^2-1)/12=21/4$,
 and the probability of exceeding $9/2$ is $4/8=1/2$.
 \item The indicator has mean $.2$ and variance $.2(.8)=.16$.
 Ten independent repetitions have a count of mean two and variance $1.6$.
 Recording one indicator ten times gives $10I$, still of mean two but
 variance $100(.16)=16$. Its only possible values are zero and ten.
 \item Without replacement, count each possible number of red balls:
 $\binom22/\binom52=1/10$ for zero reds,
 $\binom31\binom21/\binom52=6/10$ for one, and
 $\binom32/\binom52=3/10$ for two.
 The mean is $6/5$ and variance $9/25$.
 With replacement the masses are $4/25,12/25,9/25$,
 with mean $6/5$ and variance $12/25$.
 \item Examples: a fair die fits discrete uniform, whereas a loaded
 die generally does not. A success indicator fits Bernoulli, whereas
 a count taking values $0,1,2$ does not.
 A count of independent fair tosses fits binomial; repeated copies of
 one toss do not. A uniform subset of a fixed urn fits hypergeometric;
 draws with replacement generally do not.
 Independent repeated trials with a fixed success chance fit geometric;
 changing that chance after each failure generally does not.
 For a Poisson example, take the limiting count of $n$ independent
 opportunities each of probability $\lambda/n$.
 A binomial count with a fixed finite maximum and positive mean cannot
 be exactly Poisson, which has an unbounded set of possible counts.
 \item The roll count is geometric with $p=1/6$, so its mean is six,
 variance is thirty, and $\P(T>5)=(5/6)^5\approx.40188$.
 The failure count $T-1$ has mean five.
 \item $\P(X\ge2)=1-e^{-3}(1+3)=1-4e^{-3}\approx.80085$.
 Also $\E[X(X-1)]=9$, hence $\E[X^2]=9+3=12$ and variance $12-9=3$.
 \item With replacement, the count is $\Bin(4,3/10)$ with
 mean $6/5$ and variance $21/25$.
 Without replacement, the mean is still $6/5$ but the variance is
 $(21/25)(16/19)=336/475$.
 The draws have the same marginal marked probability; without
 replacement, their dependence reduces the variance.
 \item Decompose the collection time into independent geometric stages
 with success probabilities $1,5/6,4/6,3/6,2/6,1/6$.
 The mean is $6H_6=147/10=14.7$ and the variance is
 $36\sum_{j=1}^6j^{-2}-6H_6=3899/100=38.99$.
 When two faces remain, the expected further wait is $6/2+6/1=9$ rolls.
 \item Stop the stage sum after $W_{r-1}$:
 \[
  \E[T_{n,r}]=\sum_{k=0}^{r-1}\frac{n}{n-k}
   =n(H_n-H_{n-r}).
 \]
 The remaining stages have expected total $nH_{n-r}$, even conditional
 on the particular history up to seeing $r$ types. For $r=n$ this is
 zero, using $H_0=0$.
 \item Require $100e^{-m/100}\le.05$, or
 $m\ge100\log2000\approx760.09025$.
 Thus 761 draws suffice. The mean $100H_{100}\approx518.74$ describes
 an average collection time; it does not specify a $95\%$ completion
 deadline. The exponential bound is sufficient and may be conservative.
 \item With a Poisson total, the six face counts are independent
 Poisson variables of mean two. Each is zero with probability $e^{-2}$,
 so the unseen-face count is $\Bin(6,e^{-2})$, with mean $6e^{-2}$.
 The completion probability is $(1-e^{-2})^6\approx.41791$.
 For exactly 12 rolls, inclusion--exclusion instead gives
 \[
  \sum_{j=0}^6(-1)^j\binom6j(1-j/6)^{12}\approx.43782.
 \]
 Those face counts are dependent because they sum to 12. Randomizing
 the total changes the experiment and its completion probability.
 \item The median solves $e^{-\lambda m}=1/2$, giving
 $m=(\log2)/\lambda$.
 \item The standard deviation is three, so standardization gives
 $\P(8<X<13)=\Phi(1)-\Phi(-2/3)$.
 \item For $V$ uniform on $[0,1]$,
 $\E[V]=1/2$ and $\E[V^2]=1/3$, so $\Var(V)=1/12$.
 Write $X=a+(b-a)V$ to obtain mean $(a+b)/2$ and variance $(b-a)^2/12$.
\end{enumerate}

\section*{Chapter 6}
For $X$ uniform on $\{-1,0,1\}$ and $Y=X^2$, symmetry gives $\E X=\E[XY]=0$, hence covariance zero, but knowing $X$ determines $Y$. For the tower property, take the weighted average of group means; each observation receives exactly its population weight.

\section*{Chapter 7}
If $\sigma=12,n=64$, then $\operatorname{SE}(\bar X)=12/8=1.5$. The weak law follows from $\Var(\bar X_n)=\sigma^2/n$ and Chebyshev.

\section*{Chapter 8}
For 330 successes in 600 trials, $\hat p=.55$ and the plug-in standard error is
\[
\sqrt{.55(.45)/600}\approx .0203.
\]
The approximate $95\%$ interval is $.55\pm1.96(.0203)$, or about $(.510,.590)$.

\section*{Chapter 9}
Beta$(2,2)$ updated with 8 successes and 2 failures becomes Beta$(10,4)$, with posterior mean $10/14\approx.714$. The posterior predictive success probability equals the posterior mean $\alpha'/ (\alpha'+\beta')$.

\clearpage
\section*{Chapter 10}
The numbering here matches Section~\ref{sec:df-exercises}.
\begin{enumerate}
  \item The posterior probability of bias $3/4$ is
  $3^3/(3^3+3^1)=9/10$. The predictive chance of heads is
  $(1/4)(1/10)+(3/4)(9/10)=7/10$. Reordering the observed string does
  not change either calculation.
  \item Let the newly chosen biases be independent variables
  $\Theta_1,\ldots,\Theta_n$. For any specified bits,
  \[
  \P(X_1=x_1,\ldots,X_n=x_n)
  =\E\prod_{i=1}^n\Theta_i^{x_i}(1-\Theta_i)^{1-x_i}
  =\prod_{i=1}^n\frac12=2^{-n}.
  \]
  Here the tosses are generated independently given all the biases.
  The expectation factors because the biases are independent; it does
  not factor this way when a single shared $\Theta$ is reused.
  \item The ordered path has probability
  $4!1!/6!=1/30$. There are $\binom54=5$ such paths, so exactly four
  reds has probability $1/6$. After the path, the urn contains five
  red and two blue balls; the next-red probability is $5/7$.
  \item Initially $\Theta\sim\operatorname{Beta}(2,3)$. The posterior
  is Beta$(6,5)$ and the predictive probability is $6/11$.
  The updated urn also contains six red and five blue balls.
  \item Write $Y=S_n/n$ and $m=\E\Theta$. Expand
  $Y-m=(Y-\Theta)+(\Theta-m)$. The cross term has expectation zero
  because $\E[Y-\Theta\mid\Theta]=0$. Therefore
  \[
  \Var(Y)=\E[\Var(Y\mid\Theta)]+\Var(\Theta)
  =\frac{\E[\Theta(1-\Theta)]}{n}+\Var(\Theta).
  \]
  Under uniform mixing, $\E\Theta=1/2$, $\E\Theta^2=1/3$,
  so the two terms are $1/(6n)$ and $1/12$. This gives
  $\Var(S_{10}/10)=1/10$ and $\Var(S_{100}/100)=17/200$.
  The expression $1/(4n)$ wrongly treats the trials as independent
  Bernoulli$(1/2)$ variables; it misses the between-run variance.
  \item Pairwise independence implies $\Cov(X_1,X_2)=0$, hence
  $\Var(\Theta)=0$ by the representation theorem. Thus
  $\Theta=p_0$ almost surely for some constant $p_0$, and
  \eqref{eq:df} is the i.i.d.\ Bernoulli$(p_0)$ product law.
  Infinite exchangeability is the additional hypothesis. An arbitrary
  finite collection of pairwise independent events need not have it.
  \item Nonnegativity of variance gives
  \[
  0\le\Var\left(\sum_{i=1}^N X_i\right)=Nv+N(N-1)c.
  \]
  Divide by $N(N-1)$ to obtain the bound. The always-disagreeing
  pair has $v=1/4,c=-1/4$, so an exchangeable extension would require
  $-1/4\ge-1/[4(N-1)]$, possible only for $N\le2$.
  \item Fair-toss frequencies cluster near $1/2$; shared-coin
  frequencies cluster near $1/4$ and $3/4$; urn frequencies spread
  across $[0,1]$. Finite simulations fluctuate, and 100 runs will not
  give a perfectly flat histogram. The theorem identifies the limiting
  laws as $\delta_{1/2}$, $\tfrac12\delta_{1/4}+\tfrac12\delta_{3/4}$,
  and the uniform law, respectively.
  \item Conditional on the total number of ones, compare sampling
  without replacement with independent sampling at that population's
  success proportion; then average over the total. For any event,
  the absolute probability error is at most
  $\binom52/1000=.01$. With 100 draws the union bound is
  $\binom{100}{2}/1000=4.95$; capped at one, it is only the trivial
  bound. This does not say the approximation error actually equals
  one. The collision argument is useful when $n^2/N$ is small.
\end{enumerate}
